% !TEX root = ../projekt.tex
% =============================================================================
% STAV: technický obsah hotový podľa simulátora, chýbajú citácie (bod 2 zadania).
% =============================================================================
\chapter{Mechanizmus RSS}
\label{chap:rss}

RSS rozhoduje o~kanáli každého paketu v~štyroch krokoch:

\begin{enumerate}
  \item z~hlavičky paketu zostaví vstup hashu (\emph{tuple}),
  \item voliteľne ho upraví tak, aby oba smery toku dali rovnaký výsledok (\emph{symetria}),
  \item vypočíta 32-bitový hash s~tajným kľúčom,
  \item spodnými bitmi hashu vyberie položku tabuľky RETA, ktorá určí DMA kanál.
\end{enumerate}

\TODO{Obrázok tejto reťaze. Citácia pôvodnej špecifikácie RSS.}

\section{Vstup hashu}
\label{sec:tuple}

Simulátor používa 36-bajtový vstup. IPv4 adresy sú uložené ako IPv4-mapované IPv6 adresy
(\texttt{::ffff:a.b.c.d}), takže rovnaký formát pokrýva obe verzie IP.

\begin{table}[ht]
  \centering
  \small
  \caption{Rozloženie 36-bajtového vstupu hashu.}
  \label{tab:tuple}
  \begin{tabular}{r l @{\qquad} r l}
    \toprule
    \multicolumn{2}{c}{Plný tuple} & \multicolumn{2}{c}{Iba IP} \\
    Bajty & Pole & Bajty & Pole \\
    \midrule
    0--15  & zdrojová IP   & 0--15  & zdrojová IP \\
    16--17 & zdrojový port & 16--31 & cieľová IP \\
    18--33 & cieľová IP    & 32--35 & nuly \\
    34--35 & cieľový port  &        & \\
    \bottomrule
  \end{tabular}
\end{table}

Protokol do hashu nevstupuje. Porty sú zapísané v~sieťovom poradí bajtov.

\section{Symetria}
\label{sec:symetria}

Bez symetrie majú dva smery toho istého spojenia rôzny tuple a môžu skončiť v~rôznych
kanáloch. Aplikácie, ktoré spracúvajú celé spojenie, však potrebujú oba smery spolu.
Simulátor porovnáva tri režimy:

\begin{description}
  \item[none] tuple sa nemení,
  \item[xor] obe polovice sa nahradia hodnotou $\mathit{src} \oplus \mathit{dst}$,
  \item[sort] väčšia polovica (porovnanie po bajtoch) ide vždy prvá.
\end{description}

Hranica polovíc je na offsete 18 pri plnom tuple a 16 pri variante iba s~IP.

\section{Hashovacie funkcie}
\label{sec:hashe}

Referenčnou funkciou je Toeplitzov hash, štandard pre RSS. Porovnávali sa s~ním
jhash (lookup3), Chaskey, HalfSipHash a jednoduchá tabelácia (\emph{simple tabulation}).
Simulátor má navyše obálky pre multiply-shift (NH), SipHash a XOR hash.
Každá funkcia sa hodnotí so 16 náhodnými 40-bajtovými kľúčmi s~pevným semienkom.

\TODO{Ku každej funkcii 1 odsek: princíp, citácia, náročnosť v~HW. Skrátiť oproti pôvodnému
zadaniu, hash už nie je hlavná téma.}

\section{Tabuľka RETA}
\label{sec:reta}

RETA (\emph{redirection table}) je tabuľka, ktorá prevádza hash na kanál. Spodných
7 bitov hashu vyberie jednu zo 128 položiek. Položka~$i$ je predvolene nastavená na kanál
\begin{equation}
  \mathit{kanal}(i) = i \bmod n .
\end{equation}

\paragraph{Kvantizácia RETA.}
Ak $n$ nedelí 128, položky sa nedajú rozdeliť rovnomerne. Pri $n=20$ má 8 kanálov po
7 položkách a 12 kanálov po 6. Pri $n=40$ má 8 kanálov po 4 položkách a 32 kanálov po 3.
Aj dokonalý hash preto dáva nenulové preťaženie (definícia v~sekcii~\ref{sec:metriky}):
\begin{align}
  n=20: &\quad 8 \cdot \left(\tfrac{7}{128} - \tfrac{1}{20}\right) \cdot 1000 = 37{,}5\,\text{\textperthousand}, \\
  n=40: &\quad 8 \cdot \left(\tfrac{4}{128} - \tfrac{1}{40}\right) \cdot 1000 = 50{,}0\,\text{\textperthousand}.
\end{align}
Namerané hodnoty pri váhe \texttt{flow} sú 38,2 a 51,3\,\textperthousand, teda takmer celé
preťaženie pri týchto počtoch kanálov spôsobuje tabuľka, nie hash.

\section{Hlavné obmedzenie RSS}

RSS rozhoduje iba podľa hlavičky. Nevie, aký veľký tok je ani ako sú kanály práve
zaťažené. Keďže celý tok ide do jedného kanála, jeden veľký tok (\emph{elephant flow})
môže kanál preťažiť bez ohľadu na kvalitu hashu. Túto hypotézu overuje
kapitola~\ref{chap:analyza}.
